\begin{frame}{Mapa de Karnaugh: Código Gray}
    \small

    \begin{block}{Por que a ordem é diferente?}
        No Mapa de Karnaugh, células vizinhas devem diferir em
        apenas \textbf{uma variável}.
    \end{block}

    \[
        \boxed{00 \rightarrow 01 \rightarrow 11 \rightarrow 10}
    \]

    \begin{center}
        \begin{tabular}{c|c|c}
            Transição & Bit que muda & Quantidade \\ \hline
            $00 \rightarrow 01$ & $C$ & 1 \\
            $01 \rightarrow 11$ & $B$ & 1 \\
            $11 \rightarrow 10$ & $C$ & 1
        \end{tabular}
    \end{center}

    \vspace{0.3cm}

    \begin{alertblock}{Importante}
        Essa organização é chamada de \textbf{Código Gray}.

        \[
            \boxed{\text{Entre células vizinhas, apenas um bit muda.}}
        \]
    \end{alertblock}

\end{frame}